\documentclass[11pt,a4paper]{article}
\usepackage[margin=2.2cm]{geometry}
\usepackage{amsmath,amssymb,xcolor,enumitem,booktabs,graphicx,float}
\usepackage[normalem]{ulem}
\usepackage[hidelinks]{hyperref}
\setlength{\parskip}{0.6em}\setlength{\parindent}{0pt}
\definecolor{newc}{HTML}{1f6f43}
\newcommand{\new}[1]{\textcolor{newc}{#1}}
\newcommand{\old}[1]{\textcolor{gray}{\sout{#1}}}
\newcommand{\tbd}{\textcolor{orange}{TBD}}
\title{Proposed change, keeping up, the threshold and how the results are compared}
\date{23 November 2026, version 4}
\begin{document}\sloppy
\maketitle

\section*{How to read this}
Part~1 says what changed from versions 1 and 2. Part~2 explains the problem in
plain words, with the numbers behind it. Part~3 is the proposed text, place
by place. Black text stays, \new{green} is new and \old{struck} text goes.
Part~4 lists what changes outside the text, Part~5 the GAPS entries and
Part~6 the questions. Nothing below is in the thesis yet.

All numbers are 20 test seeds, counted with \texttt{served\_counts.py}
([A96]). At 16 agents the long-trained runs are used where they are
evaluated. A seed \emph{locks} when two agents hold the same task for 500
steps or more.

\section*{Part 1. What changes}
\textbf{From version 3} (author, 23 November). No controller knows which
tasks make up an order, so the Discussion says so, and deadlines on orders
become future work (3.13, [A113]).

\textbf{From version 2} (author, 23 November).
\begin{itemize}[leftmargin=1.2em]
\item \textbf{The chapter reports 16 agents, and the learned controllers at
  8 agents move to Appendix~L}, with the reason (3.12). The 16-agent policies
  were trained again for much longer, so the two fleets are not trained
  alike. Token Passing needs no training, so its 8-agent gain stays in RQ1
  (Question~1).
\item \textbf{\texttt{tab:exp:served} is easier to read.} Served share and
  seeds that keep up sit in two blocks, with no brackets (3.4).
\item \textbf{The Discussion no longer says the policies trained on
  2000-step episodes} (3.11). It says what the longer training did.
\item The lower load rates are queued for a Token Passing rerun after the
  evaluations. The strip plot goes in the chapter, and the paired
  differences stay in the text only.
\end{itemize}

\textbf{From version 1.}
\begin{itemize}[leftmargin=1.2em]
\item \textbf{The threshold goes from 0.8 to 0.9} ([A111]). At 0.8, seeds
  that pass still have a queue that grows all run long, so ``passing means
  steady'' was false for the learned controllers. At 0.9 it holds.
\item \textbf{The shipment count is named as a limit} ([A112]), with its
  size. It is measured exactly after all evaluations (queued in
  \texttt{SESSION\_CONTEXT.md}).
\item \textbf{The comparison rule stays as the thesis has it.} Version~1
  suggested comparing service time on the seeds where both runs keep up.
  That would compare only the easy seeds, since the same seeds are hard for
  every controller. So a pair is compared on service time only when both
  runs keep up on every seed, as \emph{Scoring a condition} already says.
  Every other pair is compared on served share.
\item The strip plot is redrawn at 0.9, and the counts in
  \texttt{tab:exp:served} are recounted at 0.9 for both fleets.
\end{itemize}

\section*{Part 2. The problem, in plain words}

\textbf{Two measures, each useful on one side.} Orders arrive at random, as
in Ma et al.\ (2017). A fleet that keeps up finishes what arrives, so its
tasks per step only repeat the arrival rate and cannot rank it. Its service
time can. A fleet that falls behind has a growing queue. Its service time
then grows with the length of the run, and it averages only the orders that
finished, so a worse fleet can look faster. The decentralised controller
under demand storage serves 0.53 of its work, yet its finished orders take
2255 steps, fewer than the 3355 of the centralised controller's failing
seeds, which serve more. For such a fleet the served share is the measure.
It is the work finished over the work that arrived.

\textbf{The gate.} The keep-up check decides which measure a run gets. It
is a gate, not a result. Service time is the result, and it is compared only
for runs that pass.

\textbf{Why 0.9.} Read the served share as ``of every 10 tasks that arrive,
how many get done''. A fleet that keeps up gets about 10 of 10 done.
Token Passing at the main load does, and its worst seed still serves 0.95.

\begin{table}[H]\centering\small
\begin{tabular}{lcccc}
\toprule
Token Passing, test seeds & served, mean & lowest seed & passes 0.9 & queue growth per 100 steps\\
\midrule
8 agents, every rule & 1.00 & 0.947 & 20 of 20 & $-0.02$ to 0.00\\
16 agents, every rule & 1.00 & 0.954 & 20 of 20 & $-0.05$ to 0.04\\
\bottomrule
\end{tabular}
\end{table}

At 0.8 a seed may leave 2 tasks in 10 undone and still pass. Those tasks
stay in the queue. At 0.9 the queue of a passing seed stays about flat.

\begin{table}[H]\centering\small
\begin{tabular}{lcccc}
\toprule
16 agents, long runs & passes 0.8 & queue growth & passes 0.9 & queue growth\\
\midrule
centralised, congestion-aware & 10 & $+101$ tasks & 5 & $-20$ tasks\\
section-based, congestion-aware & 9 & $+68$ & 6 & $+2$\\
section-based, fixed & 9 & $+62$ & 6 & $-14$\\
decentralised, congestion-aware & 8 & $+38$ & 6 & $-28$\\
with messages, fixed & 4 & $+75$ & 3 & $+14$\\
\bottomrule
\end{tabular}
\end{table}

The growth is over the 5000 measured steps, averaged over the passing seeds.
On the seeds that pass at 0.9 the queue is steady but long. It built up
before the measured steps, and orders take 850 to 1250 steps, against 210
to 260 under Token Passing.

\textbf{Why 0.9 costs no rerun.}
\begin{itemize}[leftmargin=1.2em]
\item Every row stores its pick and put-away shares, so keep-up is recounted
  at any threshold.
\item The checkpoint choice ([A60]) uses the check on the validation seeds.
  Recomputed at 0.9 from the saved validation rows, all seven evaluated long
  runs choose the same checkpoint.
\item At 0.9 the load sweep, counted with the served share, gives Appendix
  E's capacities (0.15, 0.3, 0.5 and 1.0), so the main loads stay 0.6 of
  capacity. At 0.8 the capacities move one step up and the main loads sit at
  about 0.5. The lower sweep rates are not yet checked with the served
  share (Question~1).
\end{itemize}

\textbf{The calculation is right.} The code (\texttt{served\_counts} in
\texttt{evaluator.py}) computes eq.~(throughput) as written. The served
share and the queue agree. Over 140 long-run seeds, the arriving minus the
finished work matches the measured queue growth with a correlation of 0.986.

\textbf{One counting limit, shipments.} An order of 5 lines needs 6 trips.
5 picks bring the units to the workstation, then 1 shipment takes the order
to a dock. The shipment is created only once all 5 picks are done. A fleet
that falls behind has many orders with some picks done, and their shipments
are not yet counted as owed. So its arriving work is counted a little low,
and its share reads a little high. Estimated from the rows, by about 0.01
for the best learned runs and about 0.05 for the weakest. Token Passing
finishes nearly every order, so its result is not affected.

\textbf{Why the count of passing seeds does not follow the mean} (from
version~1). The mean says how much is served. The count says how many seeds
land above the line, so it follows how the seeds spread. The same seeds are
hard for every controller (correlation between runs 0.57 to 0.93), and a
seed that locks drops far below the rest. Compared seed by seed, the three
controllers under congestion-aware storage cannot be told apart. The
centralised controller serves 0.020 more than the section-based one
[$-0.033$, 0.073] and 0.013 more than the decentralised one
[$-0.022$, 0.048].

\section*{Part 3. The proposed text}

\subsection*{3.1 Problem Formulation, ``Keeping up'', after eq.~(throughput)}

where the threshold \old{$f_{\mathrm{up}} \le 1$ leaves room for tasks
still under way at the horizon.}\new{$f_{\mathrm{up}}$ sits just below 1.
A fleet that keeps up serves about all of its work, and the room below 1
keeps such a fleet from failing on an unlucky seed.} The throughput
$\Lambda_T$ counts the same finished tasks, \dots

\subsection*{3.2 Experiments, ``Scoring a condition''}

\dots Only pairs in which both runs keep up under every seed are compared
on service time, and the others are reported as falling behind.
\new{A run that falls behind is compared on its served share instead, seed
by seed. Its service time grows with the length of the run, and it averages
only the orders that finished, so it would show the queue and not the
storage rule.} For the pairs that remain, \dots

\subsection*{3.3 Experiments, ``Keeping up'' (Results)}

Token Passing serves all the work \old{at both fleets}\new{at 16 agents},
under every rule (Table \texttt{tab:exp:served}).

\old{The learned controllers serve most of the work at 8 agents. They serve
between 0.81 and 0.91 of it, and each keeps up on 11 to 15 of the 20 seeds.
The decentralised controller with messages under fixed storage does best,
at 0.96 and 18 seeds.}

\old{At 16 agents they fall behind.}\new{The learned controllers fall
behind.} \old{The section-based controller serves the
most, between 0.54 and 0.69. The centralised controller serves about half.
The decentralised controller serves 0.12 under fixed storage and 0.66 under
congestion-aware storage. Only 2 to 7 seeds keep up, all of them for the
section-based or the decentralised controller.}\new{Under congestion-aware
storage the centralised, section-based and decentralised controllers each
serve about 0.8 of the work and keep up on 5 or 6 seeds. Under fixed and
demand storage they serve less, and only the section-based controller and
the one with messages keep up on any seed.} \tbd{} (follows the long runs
once all are evaluated)

\old{Passing the check does not mean serving everything. A seed passes when
it serves at least $f_{\mathrm{up}} = 0.8$ of its work. The rest waits, and
the wait grows through the run. So even on seeds that pass, a learned
controller's orders take 300 to 3200 steps, against 150 to 280 under Token
Passing. Their service times say more about the backlog than about storage,
so the learned controllers are compared by how much they serve.}

\new{A seed that keeps up is steady, but its queue can still be long. On
the learned controllers' passing seeds at 16 agents the queue no longer
grows, but it built up before the measured steps. Their orders take 850 to
1250 steps, against 210 to 260 under Token Passing. Service time is compared
only when both runs keep up on every seed, so the learned controllers are
compared by how much they serve.}

\new{\textbf{Why the count does not follow the mean.} A run can serve less
on average and still keep up on more seeds. The mean says how much a
controller serves. The count says how many seeds land above
$f_{\mathrm{up}}$, so it follows how the seeds spread
(Figure~\ref{fig:exp:seeds}).}

\new{Some seeds are harder than others for every controller. A seed's share
under one controller predicts its share under another, with a correlation
between 0.57 and 0.93. On the hardest seeds many order lines still wait for
stock at the end. Token Passing serves these seeds too, so they are harder,
not impossible.}

\new{Locks spread the seeds further. When two agents hold one task for
hundreds of steps, the work behind them waits, and the seed falls far below
the rest. The section-based controller under fixed storage locks on five
seeds, and four of them serve only 0.39 to 0.53.}

\new{So the learned controllers are compared on served share, seed by seed.
Two controllers play the same 20 seeds, so the difference on each seed
leaves out how hard the seed is. Under congestion-aware storage the
centralised controller serves 0.020 more than the section-based one
[$-0.033$, 0.073] and 0.013 more than the decentralised one
[$-0.022$, 0.048], so the three cannot be told apart. The count of seeds
that keep up stays a check and is not used to rank.}

\begin{figure}[H]\centering
\includegraphics[width=\linewidth]{../../../scripts/Images/fleet16_seeds.png}
\caption{\new{Served share of each test seed at 16 agents. A seed that locks
can fall far below the rest.}}
\label{fig:exp:seeds}
\end{figure}

\subsection*{3.4 Table \texttt{tab:exp:served}, 16 agents, two blocks}

The 8- and 32-agent columns leave the table (3.12, Question~2). The served
share and the seeds that keep up get a block each, so no number sits in
brackets.

\begin{table}[H]\centering\small
\new{Share of the arriving work served at 16 agents, mean over 20
test seeds, and the seeds that keep up.}\par\smallskip
\begin{tabular}{l ccc c ccc}
\toprule
& \multicolumn{3}{c}{served share} && \multicolumn{3}{c}{seeds that keep up}\\
\cmidrule(lr){2-4}\cmidrule(lr){6-8}
controller & fix. & dem. & cng. && fix. & dem. & cng.\\
\midrule
Token Passing & 1.00 & 1.00 & 1.00 && 20 & 20 & 20\\
centralised & 0.54$^\ast$ & 0.52$^\ast$ & 0.81 && 0$^\ast$ & 0$^\ast$ & 5\\
section-based & 0.76 & 0.54$^\ast$ & 0.79 && 6 & 0$^\ast$ & 6\\
decentralised & 0.29 & 0.53 & 0.80 && 0 & 0 & 6\\
with messages & 0.63 & 0.28$^\ast$ & 0.34$^\ast$ && 3 & 0$^\ast$ & 0$^\ast$\\
\bottomrule
\end{tabular}
\end{table}

$^\ast$ Only in this review, a first-round number until its long run is
evaluated. In the thesis every learned row is a long run, so no mark is
needed. Token Passing's shares are the test-seed replays, which also closes
the 8-agent TODO above the table.

\subsection*{3.5 Experiments, Discussion, ``The load''}

The main load assumed that every controller could keep up. It was set from
what Token Passing can serve, at \old{about} 0.6 of its capacity, before
any learned controller was trained (Appendix~E). \old{Under the served
check of the final runs, Token Passing's capacity is one tested step
higher, so the main load sits at about 0.5 of it. Every run was trained and
tested at this load, so it stays.} The assumption held for Token Passing\old{
and nearly held for the learned controllers at 8 agents. At 16 agents it
did not hold.}\new{, and at 16 agents it did not hold for the learned
controllers.}

\subsection*{3.6 Experiments, Discussion, ``Threats to validity'', new item}

\new{A shipment is created only once the last unit of its order reaches
the workstation. A fleet that falls behind therefore owes shipments that
are not counted yet, so its served share reads slightly high. At 16 agents
this is about 0.01 for the controllers that serve the most and about 0.05
for the weakest, too little to change a ranking. Token Passing finishes
nearly every order, so the main result is not affected.} \tbd{} (measured
range after the queued replay)

\subsection*{3.7 Appendix E, ``The keep-up threshold''}

The keep-up threshold is $f_{\mathrm{up}} =
\old{0.8}\new{0.9}$. \old{It leaves room for tasks that are still in
progress when the finite run ends.}\new{A fleet that keeps up finishes
about as much work as arrives, so its served share sits near 1. Under Token
Passing at the main load it is 1.00 on average at both 8 and 16 agents, and
no seed falls below 0.94. The threshold sits below that, so a fleet that
keeps up does not fail by chance. It sits high enough that a passing run
lets at most one task in ten pile up. A lower value such as 0.8 lets runs
pass whose queue grows through the whole run.} A run passes only when it
finishes at least this share of both the order work and the put-aways
released in the measured steps, counting work released earlier too.

The threshold is not varied. The backlog growth is reported next to the
check, because it shows whether unfinished work keeps growing at a rate
that passes. At every main load of the updated run the backlog shrinks
slightly. \new{With 0.9 and the served counts, the load sweep gives the
capacities of Table \texttt{tab:app:capacity}.}

In the Conclusion: ``The share 0.6 and the threshold
$f_{\mathrm{up}}=\old{0.8}\new{0.9}$ are settings of this study \dots''

\subsection*{3.8 Appendix B, Table of parameters}

$f_{\mathrm{up}}$ \dots\ \old{0.8}\new{0.9}, tested, \texttt{app:cal:load}.

\subsection*{3.9 Experiments, the 8-agent learned results leave the chapter}

\textbf{Results, opening paragraph.}
The results are for \old{8 and 16 agents, the fleets every controller could
be trained on in the time available}\new{16 agents, the fleet whose
controllers were trained in full in the time available}. Demand storage
makes Token Passing's orders 14\,\% faster at 8 agents and 20\,\% faster at
16. The learned controllers \old{keep up on most seeds at 8 agents and}
fall behind at 16, so for them the gain cannot be measured. \new{Their
results at 8 agents, from the first and shorter training, are in
Appendix~L.}

\textbf{RQ1, the controller.} The expected order, with the centralised
controller first, does not hold. \old{At 8 agents the three learned
controllers serve about the same.} At 16 agents \dots\ \tbd{} (the rest
follows the long runs, now the three are level under congestion-aware
storage).

\textbf{RQ2, congestion-aware storage}, last sentence. \old{At 8 agents the
effect is mixed.}

\textbf{RQ4, communication.} Messages help the decentralised controller
under fixed storage\old{, at both fleets}. \old{At 8 agents its served share
rises from 0.84 to 0.96, the highest of any learned run, and at 16 agents
from 0.12 to 0.37.}\new{At 16 agents its served share rises from 0.29 to
0.63.} Under demand storage messages lower the share\old{ at both fleets}.
Under congestion-aware storage they lower it\old{ at 16 agents and leave it
unchanged at 8}. \tbd{} (both follow the long runs with messages)

\textbf{Discussion, ``What was expected''.} The centralised controller did
not gain the most. \old{It never kept up at 16 agents, and at 8 agents it
served about as much as the others.}\new{At 16 agents it serves about as
much as the others under congestion-aware storage.}

\subsection*{3.10 Table \texttt{tab:exp:locks}, 16 agents only}

\begin{table}[H]\centering\small
Test seeds that lock, out of 20, at 16 agents.\par\smallskip
\begin{tabular}{l ccc}
\toprule
controller & fix. & dem. & cng.\\
\midrule
Token Passing & 0 & 0 & 0\\
centralised & 5$^\ast$ & 3$^\ast$ & 0\\
section-based & 5 & 8$^\ast$ & 3\\
decentralised & 20 & 13 & 3\\
with messages & 15 & 20$^\ast$ & 18$^\ast$\\
\bottomrule
\end{tabular}
\end{table}

The text before it: \dots\ \old{Locks are rare at 8 agents.} At 16 agents
the decentralised controller locks on \old{every seed}\new{13 to 20 seeds}
under fixed and demand storage, \old{with or without messages}\new{and with
messages on 15 under fixed storage}. Token Passing never locks. \tbd{} (the
starred cells)

\subsection*{3.11 Discussion, the training length}

\textbf{Why the learned controllers fall behind.} The learned controllers
fall behind mostly through locks\old{ and through runs that were too
short}. Two agents that each wait for the other hold one task for the rest
of the run, and the agents behind them queue. \dots\ The section-based
controller does much the same and locks less. \old{The policies also
trained and were chosen on runs of 2000 steps, while a test lasts 10000. A
policy that falls behind slowly looks fine in the short run.}\new{The first
policies also trained and were chosen on runs of 2000 steps, while a test
lasts 10000. A policy that fell behind slowly looked fine in training. The
16-agent policies were then trained again on runs of 16\,000 steps and
chosen on runs of the same length. That raised the served share of every
run. Under congestion-aware storage the decentralised controller went from
0.66 to 0.80, and under fixed storage the section-based one from 0.63 to
0.76. They still fall behind, so short runs were part of the cause, not all
of it.}

\textbf{What could have been done differently.} \old{Four changes could
have brought the learned controllers closer to Token Passing.}\new{Longer
training episodes were the first change, and they helped. Three more could
bring the learned controllers closer to Token Passing.}
\begin{itemize}[leftmargin=1.2em]
\item \old{\emph{Longer training episodes.} The policies trained on 2000
  steps and were tested on 10000. A backlog that grows slowly hardly shows in
  2000 steps, so training never punished it. Training on longer episodes, or
  choosing the checkpoint on a long validation run, targets the failure seen
  in the test.}
\item \emph{More training \new{seeds}.} Each run had \old{84 iterations
  and} one seed. \old{More iterations, or several}\new{Several} training
  seeds per controller and rule would show whether a run is simply unlucky.
\item \emph{A way out of locks.} (unchanged)
\item \emph{Starting from Token Passing.} (unchanged)
\end{itemize}

\subsection*{3.12 Appendix L, ``8 agents'', a new opening paragraph}

\new{The learned controllers at 8 agents trained only the first budget,
500\,000 steps on episodes of 2000 steps. The 16-agent policies were trained
again for 2\,000\,000 steps on episodes of 16\,000 steps, and that raised
the served share of every 16-agent run (Appendix~J). So the 8-agent
policies are less trained than the 16-agent ones. Set side by side, the two
fleets would mix the size of the fleet with the length of training. The
chapter therefore reports 16 agents, and the 8-agent runs are given here as
a first result. Token Passing needs no training, so its gain at 8 agents
stays in the chapter.}

The 8-agent tables keep their numbers, with keep-up recounted at 0.9.
Served share (seeds that keep up at 0.9):
centralised 0.86 (6), 0.90 (8), 0.81 (5); section-based 0.87 (11),
0.81 (8), 0.91 (10); decentralised 0.84 (9), 0.87 (9), 0.82 (11); with
messages 0.96 (15), 0.66 (5), 0.82 (10), for fixed, demand and
congestion-aware storage. Token Passing 1.00 (20) under every rule.

\subsection*{3.13 Discussion, orders are not seen, and deadlines as future work}

Checked in the code. The assignment rule hands out the oldest waiting task
first, to the nearest free agent (\texttt{controllers/assignment.py}), and
the network chooses only the moves. The network sees the vertex roles, the
distance to its own target, the other agents and the crowding, not the
order or the age of a task (\texttt{environment/observation.py}). The reward
pays the same for every finished task (\texttt{reward\_function.py}). Token
Passing takes the task nearest to its agent
(\texttt{planners/token\_passing.py}). Van Knippenberg et al.\ (2021) give
each job a start and a deadline, count both down in the node features, and
add a cost of 0.1 per step past a soft deadline (their Section~4.1 and
Table~1). Checked in the paper.

\textbf{New paragraph, after ``Why the learned controllers fall behind''.}

\new{\textbf{Orders are not seen.} No controller knows which tasks belong to
the same order. The assignment rule hands out the oldest task first, to the
nearest free agent, and the network then chooses only the moves. Token
Passing takes the task nearest to its agent. The reward pays the same for
every finished task. So an order with four of its five lines at the
workstation waits for its last line like any other task, although the
customer waits for the whole order. The oldest-first rule helps a little,
because the lines of an order mostly share their release time.}

\textbf{New item in ``What could have been done differently''.}
\begin{itemize}[leftmargin=1.2em]
\item \new{\emph{Deadlines on orders.} Each order could get a deadline, as
  in the time-constrained routing of van Knippenberg et al.\ (2021). Each of
  its tasks would then get its own, earlier deadline, the order's deadline
  minus the time still needed after the task, such as the shipment to the
  dock. The assignment rule could hand out the task with the least time to
  spare first, instead of the oldest. The agent could see the time left on
  its task among its features, which van Knippenberg et al.\ count down in
  the same way, and pay a small cost for every step past it. The result
  would then also count the orders that are late. Every task would carry a
  deadline, which the scope leaves out (Section~\texttt{sec:pf:scope}), so
  the problem itself would change.}
\end{itemize}

\section*{Part 4. Outside the text}
\begin{itemize}[leftmargin=1.2em]
\item \textbf{Code} (\texttt{slap-mapd-coupling}). \texttt{keep\_up\_threshold:
  0.8} becomes 0.9 in the configs. The rows already written carry
  \texttt{keeps\_up} flags at 0.8. Every number in the thesis is recounted
  from the stored shares, so the flags are not used. The author opens the PR.
\item \textbf{Scripts.} \texttt{served\_counts.py} gets \texttt{F\_UP = 0.9}.
  \texttt{fleet16\_results.py} and \texttt{appendix\_results.py} follow it,
  so Appendix L's keep-up columns are recounted.
\item \textbf{Figure.} \texttt{scripts/images/fleet16\_seeds.py}, copied to
  \texttt{thesis/Images/} once approved. Its empty columns fill as the
  remaining runs are evaluated.
\item \textbf{Queued, after every evaluation.} The order-stream replay that
  measures the shipment range (3.6).
\end{itemize}

\section*{Part 5. GAPS entries}

\textbf{[A110] The learned controllers compared on served share, seed by
seed, and the spread shown.} \emph{Where.} \texttt{5.Experiments.tex},
``Scoring a condition'', ``Keeping up'', new \texttt{fig:exp:seeds}.
\emph{What.} Service time is compared only when both runs keep up on every
seed (unchanged). Every other pair is compared on served share, with the
paired difference over the same 20 seeds and a 95\,\% interval. The count of
seeds that keep up is a gate and does not rank, because it follows the
spread of the seeds. Hard seeds are shared by every controller and locks
drop single seeds far down. Comparing only the seeds where both runs pass
was rejected, because it keeps the easy seeds. \emph{Model change.} None.

\textbf{[A111] $f_{\mathrm{up}}$ from 0.8 to 0.9.} \emph{Where.}
\texttt{3.ProblemFormulation.tex} after eq.~(throughput), Appendices B and
E, \texttt{5.Experiments.tex}, the configs. \emph{Why.} 0.8 came from the
630-step pilot of [A31], where tasks under way at the end were a large part
of the run. Under the [A96] count, tasks under way at the start and the end
roughly cancel, and a steady fleet serves about 1.0. Token Passing's lowest
test seed is 0.947 at 8 agents and 0.954 at 16. At 0.8 the learned seeds
that pass still grow their queue by 38 to 101 tasks over the measured
steps. At 0.9 they stay within $-28$ to $+14$. With the served count, 0.9
gives back Appendix E's capacities, so the main loads stay at 0.6 of
capacity. \emph{No rerun.} Keep-up is recounted from the stored shares, and
every chosen checkpoint stays the same at 0.9. \emph{Model change.} The
value only.

\textbf{[A113] Orders are not seen by any controller, and deadlines on
orders as future work.} \emph{Where.} \texttt{5.Experiments.tex},
Discussion. \emph{What.} The assignment rule, the observation and the
reward treat every task alike, and none knows the order a task belongs to,
while the result is order service time. Proposed for later work, a deadline
per order, an earlier deadline per task derived from it, assignment by the
least time to spare, the time left as an agent feature and a cost per late
step, after van Knippenberg et al.\ (2021). \emph{Model change.} None now.
The scope already excludes deadlines.

\textbf{[A112] Shipments created only when their order is complete, a
known limit.} \emph{Where.} \texttt{5.Experiments.tex}, Threats to
validity. \emph{What.} A ship task is created when the order's last unit
reaches the workstation (\texttt{\_ship}, \texttt{multi\_agent\_env.py}), so
a fleet that falls behind owes uncounted shipments and its served share
reads high. Estimated at about 0.01 for the best learned runs and 0.05 for
the weakest, measured after all evaluations by a replay of the order stream
alone. Counting one shipment per arrived order was not applied to the
finished runs, so all runs keep one count. It is the fix for later work.
\emph{Model change.} None.

\section*{Part 6. Questions}
\begin{enumerate}[leftmargin=1.5em]
\item Token Passing's 8-agent gain stays in RQ1 (\texttt{tab:exp:rq1}),
  since no training is involved. I suggest keeping it, because it shows the
  main result holds at two fleet sizes. Agreed?
\item The chapter tables had columns for 32 agents. The 32-agent policies
  were trained on the first budget too, like 8 agents. I suggest they go to
  Appendix~L with the same reason, so the chapter tables hold 16 agents
  only. Agreed?
\item Should ``Orders are not seen'' also name a first step that needs no
  deadline, an assignment rule that finishes orders already started before
  new ones? I suggest one clause at most.
\end{enumerate}

Settled in version 2: the lower load rates are queued, the strip plot goes
in the chapter, the paired differences stay in the text, and the
2000-step sentences are fixed (3.11).

\end{document}
